% Canonical LaTeX source of the rotor-chain Foster-Lyapunov benchmark draft
% (July 27, 2026), supplied verbatim by the user on 2026-08-06. Supersedes the
% PDF text extraction (2026-07-27-rotor-fl-benchmarks-original-extraction.txt)
% as the authoritative source; the reconstructed markdown transcription is now
% fully confirmed against this file.
\documentclass[11pt]{article}
\usepackage[margin=1in]{geometry}
\usepackage{amsmath,amssymb,mathtools}
\usepackage{booktabs}
\usepackage{microtype}
\usepackage{enumitem}
\usepackage{hyperref}
\usepackage[T1]{fontenc}
\usepackage{lmodern}

\newcommand{\T}{\mathbb{T}}
\newcommand{\R}{\mathbb{R}}
\newcommand{\1}{\mathbf{1}}
\newcommand{\cL}{\mathcal{L}}
\newcommand{\cO}{\mathcal{O}}
\newcommand{\ee}{\mathrm{e}}
\newcommand{\dd}{\,\mathrm{d}}
\newcommand{\wt}[1]{\widetilde{#1}}

\title{Benchmark Problems for Symbolic Foster--Lyapunov Discovery\\
       in Boundary-Thermostatted Rotor Chains}
\date{July 27, 2026}

\begin{document}
\maketitle

\section{Common setup}

\subsection{Hamiltonian rotor chain}
For $N\ge 2$, let
\[
q=(q_1,\ldots,q_N)\in\T^N,
\qquad
p=(p_1,\ldots,p_N)\in\R^N.
\]
Here $\T=\R/(2\pi\mathbb Z)$ denotes the one-dimensional torus: angles that differ by an integer multiple of $2\pi$ represent the same rotor position. Thus $\T^N$ is the Cartesian product of $N$ circles, and every $q_i$ is understood modulo $2\pi$; $p_i$ is its conjugate momentum. Define the bond variables
\[
r_i=q_{i+1}-q_i\pmod{2\pi},
\qquad i=1,\ldots,N-1,
\]
and consider the Hamiltonian
\begin{equation}
H_N(q,p)=\frac12\sum_{i=1}^{N}p_i^2+\sum_{i=1}^{N-1}W_i(r_i).
\end{equation}
The general theory allows smooth, periodic, non-degenerate interactions. To make the four benchmarks directly comparable, we fix
\begin{equation}
W_i(r)=W(r)=1-\cos r,
\qquad
w(r)=W'(r)=\sin r.
\end{equation}
Adding the constant $1$ to $-\cos r$ does not change the dynamics, but makes $H_N\ge0$.

\subsection{Endpoint Langevin baths}
For $N\ge3$, only rotors $1$ and $N$ are coupled to heat baths. We use the concrete parameters
\begin{equation}
\gamma_1=\gamma_N=1,
\qquad T_1=1,
\qquad T_N=2,
\end{equation}
and zero external torques. The stochastic dynamics are
\begin{subequations}
\begin{align}
\dd q_i &=p_i\dd t, &&i=1,\ldots,N,\\
\dd p_1&=\bigl(w(r_1)-p_1\bigr)\dd t+\sqrt2\,\dd B_t^1,\\
\dd p_i&=\bigl(w(r_i)-w(r_{i-1})\bigr)\dd t,
&&i=2,\ldots,N-1,\\
\dd p_N&=\bigl(-w(r_{N-1})-p_N\bigr)\dd t+2\,\dd B_t^N.
\end{align}
\end{subequations}
where $B^1$ and $B^N$ are independent standard Brownian motions. More generally, the noise coefficient at endpoint $b$ is $\sqrt{2\gamma_bT_b}$.

Let $F_i(q)=-\partial_{q_i}H_N(q,p)$. The infinitesimal generator acting on a $C^2$ function $f$ is
\begin{equation}
\cL_Nf=\sum_{i=1}^{N}\bigl(p_i\partial_{q_i}f+F_i(q)\partial_{p_i}f\bigr)
-p_1\partial_{p_1}f-p_N\partial_{p_N}f
+\partial_{p_1p_1}f+2\partial_{p_Np_N}f.
\end{equation}
In particular,
\begin{equation}
\cL_NH_N=(1-p_1^2)+(2-p_N^2).
\end{equation}
Thus $H_N$ alone does not detect energy trapped in the undamped interior rotors and is not a sufficient Foster--Lyapunov function.

\subsection{One-bath $N=2$ specialization}
The baby benchmark is deliberately not obtained by identifying the two endpoints in (4). Instead, only rotor $1$ is thermostatted, with $\gamma_1=T_1=1$, and rotor $2$ has neither friction nor noise:
\begin{subequations}
\begin{align}
\dd q_i&=p_i\dd t, &&i=1,2,\\
\dd p_1&=\bigl(w(r_1)-p_1\bigr)\dd t+\sqrt2\,\dd B_t^1,\\
\dd p_2&=-w(r_1)\dd t.
\end{align}
\end{subequations}
Its generator is
\begin{equation}
\cL_2f=p_1\partial_{q_1}f+p_2\partial_{q_2}f
+w(r_1)(\partial_{p_1}-\partial_{p_2})f
-p_1\partial_{p_1}f+\partial_{p_1p_1}f,
\end{equation}
and
\begin{equation}
\cL_2H_2=1-p_1^2.
\end{equation}
Thus the unthermostatted second rotor creates the same fast-rotor obstruction as an interior rotor in a longer chain.

\subsection{Target certificate}
For $R>0$, let
\[
K_R=\{(q,p):H_N(q,p)\le R\}.
\]
The model must generate an explicit, periodic-in-$q$, $C^2$ function
\[
V_N:\T^N\times\R^N\longrightarrow[1,\infty)
\]
with compact sublevel sets, together with constants $c,C,R>0$, such that
\begin{equation}
\cL_NV_N\le C\1_{K_R}-\phi_N(V_N)
\qquad\text{on }\T^N\times\R^N.
\end{equation}
We use the parametric drift family
\begin{equation}
\phi_{\alpha,c}(v)=c\frac{v}{(\alpha+1+\log v)^\alpha},
\qquad v\ge1.
\end{equation}
The offset $\alpha+1$ makes this representative increasing and concave on $[1,\infty)$; its precise value is otherwise inessential. The significant feature is the large-$v$ behavior $v/(\log v)^\alpha$.

These are stochastic Foster--Lyapunov problems. Therefore the reward and falsification code must evaluate the full generator (8) or (5), including the second derivatives in every noisy momentum (one for $N=2$ and two for $N\ge3$). Replacing $\cL_NV$ by the first-order Lie derivative would define a different problem.

\begin{center}
\begin{tabular}{ccccc}
\toprule
$N$ & state dimension & status & target $\alpha$ & associated rate\\
\midrule
2 & 4  & proved      & 1 & $\exp(-ct^{1/2})$\\
3 & 6  & proved      & 1 & $\exp(-ct^{1/2})$\\
4 & 8  & proved      & 1 & $\exp(-ct^{1/2})$\\
5 & 10 & conjectural & 3 & $\exp(-ct^{1/4})$\\
\bottomrule
\end{tabular}
\end{center}

\section{The four testing problems}

\noindent\textbf{Testing problem 1 ($N=2$: one fast unthermostatted rotor).}
\emph{Use the one-bath dynamics (7). Explicitly,}
\begin{align*}
H_2&=\frac12(p_1^2+p_2^2)+(1-\cos(q_2-q_1)),\\
\dd p_1&=\bigl(\sin(q_2-q_1)-p_1\bigr)\dd t+\sqrt2\,\dd B_t^1,\\
\dd p_2&=-\sin(q_2-q_1)\dd t.
\end{align*}
\emph{Find a globally $C^2$, coercive $V_2$ and $c,C,R>0$ such that}
\begin{equation}
\cL_2V_2\le C\1_{K_R}-c\frac{V_2}{2+\log V_2}.
\end{equation}
\emph{The output must be periodic in $(q_1,q_2)$ and coercive in $(p_1,p_2)$, including along the fast unthermostatted direction $|p_2|\to\infty$ with $p_1$ bounded.}

\medskip
\noindent\textbf{Known reference information.}
Cuneo and Poquet study precisely this two-rotor, one-bath model and prove that the stretched-exponential scale $\exp(-ct^{1/2})$ cannot be improved~[1]. They also observe that the matching upper bound follows from the three-rotor Foster--Lyapunov proof by omitting the third rotor~[2]. The resulting one-branch reference construction has, for sufficiently small $0<\theta<1$, the structure
\begin{equation}
V_2^{\mathrm{ref}}=1+\ee^{\theta H_2}
+A\rho_2(p)\wt p_2^{\;2}\exp\!\left(\frac\theta2\wt p_2^{\;2}\right).
\end{equation}
We fix the gate explicitly. Let
\begin{equation}
b(u)=
\begin{cases}
0,&u\le0,\\
\exp(-1/u),&u>0,
\end{cases}
\qquad
\chi(z)=\frac{b(z-1)}{b(z-1)+b(4-z)}.
\end{equation}
Thus $\chi(z)=0$ for $z\le1$ and $\chi(z)=1$ for $z\ge4$. Taking $k=2$ and $R_\rho=4$, define
\begin{equation}
S_2(p_1)=p_1^4+4,
\qquad
\rho_2(p_1,p_2)=\chi\!\left(\frac{p_2^2}{S_2(p_1)^2}\right).
\end{equation}
In particular, $\rho_2=0$ when $|p_2|\le p_1^4+4$, and $\rho_2=1$ when $|p_2|\ge2(p_1^4+4)$. To regularize the inverse powers appearing below, set
\begin{equation}
I_j(p_1,p_2)=
\begin{cases}
\rho_2(p_1,p_2)/p_2^j,&p_2\ne0,\\
0,&p_2=0,
\end{cases}
\qquad j=1,\ldots,4.
\end{equation}
Each $I_j$ is globally $C^\infty$, since $\rho_2$ vanishes on a neighborhood of $p_2=0$. The corrected fast momentum is obtained by averaging the oscillatory interaction force. With the regularization above, one concrete normal form for the cosine benchmark is
\begin{align}
\wt p_2={}&p_2-(\cos r_1)I_1-(p_1\cos r_1)I_2\notag\\
&-(p_1^2\cos r_1+p_1\sin r_1+\cos^2r_1)I_3\notag\\
&+\left(-p_1^3\cos r_1-3p_1^2\sin r_1-3p_1\cos^2r_1+p_1\cos r_1
-\frac12\sin r_1\cos r_1+2\sin r_1\right)I_4.
\end{align}
On the fully selected fast branch $\rho_2=1$, this agrees with the inverse-power normal form and satisfies
\begin{equation}
\dd\wt p_2=
\left(-\frac{1}{2p_2^3}+\cO_f(p_2^{-4})\right)\dd t
+\left(-\frac{\sqrt2\cos r_1}{p_2^2}+\cO_f(p_2^{-3})\right)\dd B_t^1.
\end{equation}
Here and below, $\cO_f$ denotes the weighted fast-region remainders: their coefficients may depend polynomially on the thermostatted endpoint momenta but are controlled by the corresponding gate. The singular expression (17) has become globally smooth through (16). There are constants $a_1,a_2>0$ such that
\begin{equation}
1+a_1\ee^{\theta H_2}\le V_2^{\mathrm{ref}}\le a_2(1+p_2^2)\ee^{\theta H_2}.
\end{equation}
The separate test function constructed explicitly in~[1] establishes the lower relaxation bound; it has a positive upper drift and should not be confused with the negative-drift Foster--Lyapunov reference (13).

\medskip
\noindent\textbf{Testing problem 2 ($N=3$: one fast interior rotor).}
Set $N=3$ in (1)--(4). Explicitly,
\begin{align*}
H_3&=\frac12(p_1^2+p_2^2+p_3^2)+(1-\cos(q_2-q_1))+(1-\cos(q_3-q_2)),\\
\dd p_1&=\bigl(\sin(q_2-q_1)-p_1\bigr)\dd t+\sqrt2\,\dd B_t^1,\\
\dd p_2&=\bigl(\sin(q_3-q_2)-\sin(q_2-q_1)\bigr)\dd t,\\
\dd p_3&=\bigl(-\sin(q_3-q_2)-p_3\bigr)\dd t+2\,\dd B_t^3.
\end{align*}
Find $V_3$ and $c,C,R>0$ satisfying
\begin{equation}
\cL_3V_3\le C\1_{K_R}-c\frac{V_3}{2+\log V_3}.
\end{equation}
The output must be coercive in $(p_1,p_2,p_3)$ and globally $C^2$.

\medskip
\noindent\textbf{Known reference information.}
This target is rigorous for sufficiently small exponential weight $\theta>0$~[2]. The published construction has the structure
\begin{equation}
V_3^{\mathrm{ref}}=1+\ee^{\theta H_3}
+A\rho_3(p)\wt p_2^{\;2}\exp\!\left(\frac\theta2\wt p_2^{\;2}\right),
\end{equation}
where $\rho_3$ is a smooth gate selecting the regime in which $|p_2|$ dominates the endpoint momenta, and $\wt p_2$ is a corrected fast momentum whose dynamics become slow after averaging the rapidly oscillating forces. For the cosine benchmark, the first two counterterms are explicitly
\begin{equation}
\wt p_2=p_2-\frac{\cos r_1+\cos r_2}{p_2}
-\frac{p_1\cos r_1+p_3\cos r_2}{p_2^2}+\cO_f(p_2^{-3}).
\end{equation}
The higher counterterms in (22) are selected recursively, not arbitrarily, so that on the gated fast branch
\begin{align}
\dd\wt p_2={}&\left(-\frac1{p_2^3}+\cO_f(p_2^{-4})\right)\dd t\notag\\
&+\left(-\frac{\sqrt2\cos r_1}{p_2^2}+\cO_f(p_2^{-3})\right)\dd B_t^1\notag\\
&+\left(-\frac{2\cos r_2}{p_2^2}+\cO_f(p_2^{-3})\right)\dd B_t^3.
\end{align}
As in the $N=2$ case, this fast-branch expression is smoothly extended where $\rho_3=0$. The reference theorem gives constants $a_1,a_2>0$ such that
\begin{equation}
1+a_1\ee^{\theta H_3}\le V_3^{\mathrm{ref}}\le a_2(1+p_2^2)\ee^{\theta H_3}.
\end{equation}
For a strict blind test, equations (21)--(24) should be withheld from the generator and used only for comparison.

\medskip
\noindent\textbf{Testing problem 3 ($N=4$: two fast interior rotors and resonances).}
Set $N=4$. The interior momentum equations are
\begin{align*}
\dd p_2&=\bigl(\sin(q_3-q_2)-\sin(q_2-q_1)\bigr)\dd t,\\
\dd p_3&=\bigl(\sin(q_4-q_3)-\sin(q_3-q_2)\bigr)\dd t,
\end{align*}
with the endpoint equations given by (4). Find a globally $C^2$, coercive $V_4$ and $c,C,R>0$ such that
\begin{equation}
\cL_4V_4\le C\1_{K_R}-c\frac{V_4}{2+\log V_4}.
\end{equation}
The certificate must cover all high-energy regimes: only $p_2$ fast, only $p_3$ fast, both fast and non-resonant, and both fast and resonant.

\medskip
\noindent\textbf{Known reference information.}
This target is also rigorous~[3]. For any
\[
0<\theta<\min(T_1^{-1},T_4^{-1})=\frac12,
\]
the reference construction patches several regime-specific functions:
\begin{equation}
V_4^{\mathrm{ref}}=1+\ee^{\theta H_4}+\psi_2(p)V_2+\psi_3(p)V_3+M\psi_c(p)V_c.
\end{equation}
Here $\psi_2,\psi_3,\psi_c$ are smooth regime gates. The one-fast-rotor pieces have the asymptotic form
\[
V_j\asymp\exp\!\left(|\wt p_j|^a+\frac\theta2\wt p_j^{\;2}\right),
\qquad j=2,3,\quad 0<a<1,
\]
whereas the both-fast piece has the form
\[
V_c\asymp\wt H_c\ee^{\theta\wt H_c},
\qquad
\wt H_c\asymp\frac12(p_2^2+p_3^2).
\]
The tildes denote averaging and resonance counterterms, not merely a change of notation. The resulting global function obeys
\begin{equation}
1+\ee^{\theta H_4}\le V_4^{\mathrm{ref}}
\le a_3\bigl(\ee^{|p_2|^a}+\ee^{|p_3|^a}\bigr)\ee^{\theta H_4}
\end{equation}
for some $a_3>0$. Again, the reference structure can either be supplied as a grammar prior or withheld for a genuinely blind test.

\medskip
\noindent\textbf{Testing problem 4 ($N=5$: conjectural central-rotor bottleneck).}
Set $N=5$. In addition to the endpoint equations, the interior dynamics are
\begin{align*}
\dd p_2&=\bigl(\sin(q_3-q_2)-\sin(q_2-q_1)\bigr)\dd t,\\
\dd p_3&=\bigl(\sin(q_4-q_3)-\sin(q_3-q_2)\bigr)\dd t,\\
\dd p_4&=\bigl(\sin(q_5-q_4)-\sin(q_4-q_3)\bigr)\dd t.
\end{align*}
Find or falsify the existence of a globally $C^2$, coercive $V_5$ and $c,C,R>0$ satisfying the conjectural inequality
\begin{equation}
\cL_5V_5\le C\1_{K_R}-c\frac{V_5}{(4+\log V_5)^3}.
\end{equation}
The phrase ``find or falsify'' is important: failure of a particular symbolic ansatz is not a disproof, but an explicit sequence of states along which its drift inequality fails is a useful counterexample for further training.

\medskip
\noindent\textbf{Conjecture 1 (Benchmark conjecture for $N=5$).}
For the chain above, there are $\theta,c,C,R>0$ and a smooth $V_5\ge1$ with
\[
\log V_5=\theta H_5+o(H_5)\qquad\text{as }H_5\to\infty
\]
such that (28) holds.

The exponent $3$ is inferred from the conjectured mixing law in~[3, Remark 5.3]. If $\phi(v)\asymp v/(\log v)^\alpha$, then
\[
H_\phi(u)=\int_1^u\frac{\dd s}{\phi(s)}\asymp(\log u)^{\alpha+1},
\]
and the corresponding stretched-exponential power is $1/(\alpha+1)$. The prediction
\[
k_N=\frac1{2\lceil N/2\rceil-2}
\]
gives $k_5=1/4$, hence $\alpha=3$.

One heuristic starting point for the generator is
\begin{equation}
V_5^{\mathrm{ansatz}}=1+\ee^{\theta H_5}
+A\psi_3(p)\wt p_3^{\;6}\exp\!\left(\frac\theta2\wt p_3^{\;2}\right)
+\text{corrections for the other fast and resonant regimes}.
\end{equation}
This ansatz is motivated by the unproved effective scaling
\[
\cL_5\wt p_3\sim-\kappa\wt p_3^{-7}
\quad\Longrightarrow\quad
\frac{\cL_5V_5}{V_5}\sim-\wt p_3^{-6}\sim-(\log V_5)^{-3}.
\]
Both the scaling and the required correction terms are conjectural. A serious search must also treat high energy shared among $(p_2,p_3)$, $(p_3,p_4)$, or $(p_2,p_3,p_4)$, where new resonances can arise.

\section{Suggested evaluation protocol}
For each $N$, the model receives $H_N$, $\cL_N$, the target exponent $\alpha$, and an allowed symbolic grammar. It returns an expression tree for $V_N$ and numerical or symbolic values for any free parameters. Evaluation should be split into the following stages.

\begin{enumerate}[leftmargin=*,itemsep=0.75em]
\item \textbf{Syntax and regularity.} Check that the expression is real, periodic in each angle, globally defined, and $C^2$. Denominators such as $1/p_i$ require smooth cutoffs that vanish near $p_i=0$.

\item \textbf{Positivity and coercivity.} Search for violations of $V_N\ge1$ and verify growth as $\lVert p\rVert\to\infty$.

\item \textbf{Generator falsification.} Maximize
\[
\cL_NV_N+c\frac{V_N}{(\alpha+1+\log V_N)^\alpha}
\]
outside $K_R$, with adversarial sampling concentrated in fast-rotor and near-resonant regimes.

\item \textbf{Tail certification.} A bounded-box numerical check is not a global Foster--Lyapunov proof. Split momentum space into asymptotic regimes and prove the sign of the leading terms in each regime.

\item \textbf{Blind versus guided scores.} Run a blind task using only the dynamics, then a guided task in which the structural reference forms (13), (21), (26), or (29) are included in the grammar prior. Report both scores.
\end{enumerate}

A practical first grammar should contain arithmetic operations, integer and learned real constants, $\sin$, $\cos$, $\exp$, $\log$, absolute or smooth absolute value, and a fixed smooth cutoff primitive. The $N=4$ and $N=5$ problems are unlikely to admit short global expressions without some mechanism for regime-dependent patching.

\section{Questions left open in this draft}
\begin{enumerate}[leftmargin=*,itemsep=0.75em]
\item Should the benchmark require a fully global certificate, or should it also include a finite-domain discovery track compatible with existing numerical falsifiers?
\item Should the bath temperatures and friction coefficients be fixed as above or sampled during training to test parameter generalization, separately for the one-bath $N=2$ and two-bath $N\ge3$ cases?
\item For $N=5$, should $\alpha=3$ be supplied to the model, or should the model jointly search over $\alpha$ and $V_5$?
\item Which smooth cutoff functions and piecewise constructions should be included as atomic tokens in the discrete diffusion vocabulary?
\end{enumerate}

\begin{thebibliography}{9}
\bibitem{cuneo-poquet-2017}
Noé Cuneo and Christophe Poquet.
On the relaxation rate of short chains of rotors interacting with Langevin thermostats.
\emph{Electronic Communications in Probability}, 22:1--8, 2017. Paper No.~35.
\href{https://doi.org/10.1214/17-ECP62}{doi:10.1214/17-ECP62}.

\bibitem{cuneo-eckmann-poquet-2015}
Noé Cuneo, Jean-Pierre Eckmann, and Christophe Poquet.
Non-equilibrium steady state and subgeometric ergodicity for a chain of three coupled rotors.
\emph{Nonlinearity}, 28(7):2397--2421, 2015.
\href{https://doi.org/10.1088/0951-7715/28/7/2397}{doi:10.1088/0951-7715/28/7/2397}.

\bibitem{cuneo-eckmann-2016}
Noé Cuneo and Jean-Pierre Eckmann.
Non-equilibrium steady states for chains of four rotors.
\emph{Communications in Mathematical Physics}, 345(1):185--221, 2016.
\href{https://doi.org/10.1007/s00220-015-2550-2}{doi:10.1007/s00220-015-2550-2}.

\bibitem{zou-et-al-2025}
Haohan Zou, Jie Feng, Hao Zhao, and Yuanyuan Shi.
Analytical Lyapunov function discovery: An RL-based generative approach.
In Aarti Singh, Maryam Fazel, Daniel Hsu, Simon Lacoste-Julien, Felix Berkenkamp, Tegan Maharaj, Kiri Wagstaff, and Jerry Zhu, editors,
\emph{Proceedings of the 42nd International Conference on Machine Learning}, volume 267 of \emph{Proceedings of Machine Learning Research}, pages 80776--80804. PMLR, 2025.
\url{https://proceedings.mlr.press/v267/zou25a.html}.
\end{thebibliography}

\end{document}
